% TC-0018 controlled test-case fragment.
\def\TCID{TC-0018}
\def\TCTitle{Optional GUI process startup}
\def\TCPurpose{Verify that WCRT supplies an explicitly selected GUI PE entry
point that invokes ANSI \texttt{WinMain} without affecting other startup paths.}
\def\TCPriority{\TCPriorityHigh}
\def\TCRequirementRef{REQ-0018}
\def\TCDesignRef{ADR-0001; WCRT Phase 0 process-startup objectives; no C89
draft clause.}
\def\TCFeatureRef{\texttt{wcrt-startup-gui.o}; \texttt{libwcrt.a};\par
\texttt{tests/c89/run-tc-0018.ps1}.}
\def\TCTestTechnique{Use-case testing, equivalence partitioning of GUI command
lines and show states, and binary inspection.}
\def\TCPreconditions{TinyCC, a completed WCRT static build, the optional GUI
startup object, PE inspection tooling, and controlled child-process observation
are available for the target architecture.}
\def\TCEnvironment{Windows x86, x64, and ARM64 release-matrix environments;
native execution is required for final release evidence.}
\def\TCAssumptions{The test uses the documented GUI-subsystem link recipe and
no toolchain-provided C runtime startup.}
\def\TCInputData{Command lines with no application arguments, spaces, empty
arguments, and quoted text; default and explicitly supplied initial-show state.}
\def\TCInitialState{No test executable or evidence from an earlier execution
is present in the isolated output directory.}
\def\TCProcedure{\par\begin{enumerate}
\item Verify that an ordinary \texttt{libwcrt.a} link does not select the WCRT
GUI entry point.
\item Link a \texttt{WinMain} consumer with the documented no-toolchain-runtime
command, \texttt{wcrt-startup-gui.o}, and \texttt{libwcrt.a}.
\item Inspect the PE machine type, GUI subsystem, entry point, and import table;
reject any host C runtime dependency.
\item Execute the consumer as a controlled child for each command-line and
initial-show-state partition.
\item Verify the current module instance, null previous instance, command line
without the program name, and applicable initial-show value passed to
\texttt{WinMain}.
\item Return a controlled value from \texttt{WinMain} and verify process
termination. Register an \texttt{atexit} callback and verify that it runs.
\end{enumerate}}
\def\TCExpectedResult{The artifact is a GUI-subsystem PE of the expected
architecture with no host C runtime import. \texttt{WinMain} receives every
required value, its \texttt{atexit} callback runs, and its return value becomes
the process exit status.}
\def\TCPostConditions{The child has terminated, and generated executables and
evidence remain only in the isolated output directory.}
\TCTemplate
